% Chapter 6

\chapter{Conclusion}

\label{Chapter6}

%---------------------------------------------------------------------------------------- 
\newcommand{\keyword}[1]{\textbf{#1}}
\newcommand{\tabhead}[1]{\textbf{#1}}
\newcommand{\code}[1]{\texttt{#1}}
\newcommand{\file}[1]{\texttt{\bfseries#1}}
\newcommand{\option}[1]{\texttt{\itshape#1}}
\renewcommand{\figurename}{Figure.}
\renewcommand{\tablename}{Table.}

%----------------------------------------------------------------------------------------
Completing all the objectives of this master's thesis verified the research hypotheses and led to the following conclusions:

\begin{enumerate}
  \item DFT-derived descriptors (i.e., HOMO–LUMO gap, solvent accessible surface area, Hartree–Fock energy) improve PFAS half-life classification when integrated with biological descriptors (i.e., Sex, Species) and physicochemical properties (i.e., logP);
  
  \item The choice of computational method influences descriptor sensitivity and overall model performance;
  
  \item A balanced, mid-level DFT approach (e.g., M06-2X/6-31++G(d,p)) often provides the optimal trade-off between accuracy and efficiency;
  
  \item Ensemble learning (e.g., XGBoost) effectively utilizes the descriptor set; however, considerations of interpretability and computational cost remain important.
\end{enumerate}

These findings offer a practical workflow for classifying PFAS persistence, advancing predictive toxicokinetic modeling, and supporting safer chemical design and informed regulatory decision-making.

